%% ============================================================
%%  Capítulo 4 — Implementação de JML e MFA
%% ============================================================
\chapter{Implementação de JML e MFA}
\label{chap:jml_mfa}

\section{Ciclo de Vida JML}

O ciclo de vida Joiner–Mover–Leaver (\gls{jml}) automatiza o provisionamento, a transição e a revogação de acessos.

\subsection{Arquitetura}

Os três fluxos são implementados sobre a Keycloak Admin \gls{rest} \gls{api}, garantindo que toda a gestão de identidades é centralizada no \gls{idp}. Cada fluxo está disponível como:
\begin{itemize}
  \item Script \gls{cli}: \texttt{jml/joiner.py}, \texttt{mover.py}, \texttt{leaver.py}
  \item Endpoint \gls{rest}: \texttt{POST /admin/jml/\{joiner|mover|leaver\}}
\end{itemize}

\subsection{Os Três Fluxos}

\begin{table}[h!]
\centering
\caption{Fluxos JML: ações e efeito de segurança imediato}
\label{tab:jml}
\begin{tabularx}{\textwidth}{|l|X|l|}
\hline
\textbf{Fluxo} & \textbf{Ações} & \textbf{Efeito imediato} \\ \hline
Joiner & Cria utilizador, atribui papel e grupo inicial & Acesso inicial controlado \\ \hline
Mover & Troca papel/grupo, \textbf{revoga sessões ativas} & Novo acesso após reautenticação \\ \hline
Leaver & Desativa conta (\texttt{enabled:false}), remove papéis, \textbf{revoga sessões} & Bloqueio imediato de acesso \\ \hline
\end{tabularx}
\end{table}

\begin{figure}[h!]
  \centering
  \includegraphics[width=0.95\textwidth]{fluxo_jml.png}
  \caption{Fluxos JML: criação, transição de papel/grupo, desativação e revogação de sessões.}
  \label{fig:jml}
\end{figure}

\subsection{Decisão Técnica: Revogação Explícita de Sessões}

Sem revogação explícita, um utilizador movido ou desativado manteria o acesso durante o tempo de vida do \gls{jwt} (por omissão, 300 segundos). A chamada \texttt{DELETE /admin/realms/iam-tp/users/\{id\}/sessions} força a reautenticação imediata, eliminando esta janela de acesso indevido.

\begin{lstlisting}[language=Python, caption={Revogação de sessões no mover.py},label={lst:revoke}]
# Apos trocar papel e grupo:
admin_token = get_admin_token()
user_id = get_user_id(username, admin_token)
revoke_url = f"{KEYCLOAK_URL}/admin/realms/{REALM}/users/{user_id}/sessions"
httpx.delete(revoke_url, headers={"Authorization": f"Bearer {admin_token}"})
\end{lstlisting}

\section{Autenticação Multifator (MFA/TOTP)}

\subsection{Configuração no Keycloak}

A \gls{mfa} é implementada com as seguintes características:

\begin{enumerate}
  \item O fluxo de autenticação \textit{browser} inclui a execução \texttt{auth-otp-form} (obrigatória para o papel \texttt{admin});
  
  \item O utilizador \texttt{admin.user} tem a \textit{Required Action} \texttt{CONFIGURE\_TOTP} ativa;
  
  \item No primeiro acesso via \textit{browser}, o Keycloak força a leitura do \textit{QR Code} e a configuração da aplicação TOTP (Google Authenticator, Authy, etc.);
  
  \item Após configuração bem-sucedida, o \gls{jwt} passa a incluir \texttt{"amr": ["otp"]} como prova de segundo fator.
\end{enumerate}

\subsection{Validação na API}

O endpoint \texttt{/admin/mfa-area} exige cumulativamente:
\begin{itemize}
  \item Papel \texttt{admin}
  \item \textit{Claim} \texttt{amr=otp}
\end{itemize}

Esta separação demonstra que a \gls{api} valida a prova de \gls{mfa} de forma independente do Keycloak:
\begin{itemize}
  \item Um token obtido via \gls{ropc} (sem \gls{mfa}) retorna \textbf{403}
  \item Um token obtido via \textit{browser flow} com \gls{totp} bem-sucedido retorna \textbf{200}
\end{itemize}

\subsection{Fluxo de TOTP}

\begin{enumerate}
  \item Utilizador faz login no Keycloak
  \item Sistema deteta papel \texttt{admin} e requer \gls{otp}
  \item Utilizador escaneia \textit{QR Code} com aplicação TOTP (primeira vez)
  \item Utilizador insere código de 6 dígitos
  \item Keycloak valida o código (período de 30 segundos)
  \item Token é gerado com \texttt{amr: ["otp"]}
\end{enumerate}

\section{Auditoria de Eventos}

\subsection{Tipos de Eventos Registados}

O Keycloak regista automaticamente dois tipos de eventos com persistência em PostgreSQL:

\begin{itemize}
  \item \textbf{Eventos de utilizador}: login, falha de login, logout, \textit{refresh} de token, configuração de \gls{totp} e falhas de \gls{otp};
  
  \item \textbf{Eventos de administração}: criação/remoção de utilizadores, alteração de papéis e grupos, atualização de clientes.
\end{itemize}

\subsection{Acesso à Auditoria}

O endpoint \texttt{GET /admin/audit} (protegido por papel \texttt{admin}) agrega ambos os tipos numa única resposta e suporta filtragem por tipo de evento.

\subsection{Exemplos de Eventos}

\subsubsection{Evento de Login bem-sucedido}

\begin{lstlisting}[caption={Excerto de evento de auditoria — login bem-sucedido},label={lst:audit_evento}]
{
  "id": "b3a7c1f2-8d4e-4a2b-9f1c-3e5d7a8b2c4d",
  "time": 1748430600000,
  "type": "LOGIN",
  "realmId": "iam-tp",
  "clientId": "fastapi-client",
  "userId": "a1b2c3d4-e5f6-7890-abcd-ef1234567890",
  "sessionId": "c9d8e7f6-a5b4-3c2d-1e0f-9a8b7c6d5e4f",
  "ipAddress": "172.18.0.1",
  "details": {
    "auth_method": "openid-connect",
    "auth_type": "code",
    "response_type": "code",
    "redirect_uri": "http://localhost:8000/auth/callback",
    "username": "colaborador.user"
  }
}
\end{lstlisting}

\subsubsection{Evento de Administração — Criação de Utilizador (Joiner)}

\begin{lstlisting}[caption={Excerto de evento de administração — criação de utilizador (Joiner)},label={lst:audit_admin}]
{
  "id": "f4e3d2c1-b0a9-8765-4321-fedcba987654",
  "time": 1748431200000,
  "operationType": "CREATE",
  "resourceType": "USER",
  "realmId": "iam-tp",
  "representation": "{\"username\":\"demo.user\",\"enabled\":true}",
  "resourcePath": "users/f1e2d3c4-a5b6-7890-abcd-ef0123456789",
  "authDetails": {
    "realmId": "iam-tp",
    "clientId": "admin-cli",
    "userId": "admin-service-account",
    "ipAddress": "172.18.0.1"
  }
}
\end{lstlisting}

\subsection{Evidência Consultável}

A evidência consultável completa está disponível no dashboard interativo (aba \textbf{Auditoria}), com:
\begin{itemize}
  \item Filtragem por tipo de evento (\texttt{LOGIN}, \texttt{LOGIN\_ERROR}, \texttt{UPDATE\_TOTP}, etc.)
  \item Resumo estatístico por categoria
  \item Visualização de eventos de utilizador e administração
\end{itemize}
